\chapter{CODE}
\label{app:code}

\section{Main MATLAB Code Structure}
\begin{lstlisting}[language=Matlab]
clc;
clear;
close all;

data = readtable('MSME_Energy_Data.xlsx');

production = data.Production;
grid_energy = data.Grid_Energy;
dg_energy = data.DG_Energy;

valid_idx = production > 0 & grid_energy >= 0 & dg_energy >= 0;
production = production(valid_idx);
grid_energy = grid_energy(valid_idx);
dg_energy = dg_energy(valid_idx);

total_energy = grid_energy + dg_energy;

efficiency = production ./ total_energy;
EPU = total_energy ./ production;

grid_tariff = 6;
dg_tariff = 15;
emission_factor = 0.8;

cost = grid_energy .* grid_tariff + dg_energy .* dg_tariff;
co2 = dg_energy .* emission_factor;

linear_model = polyfit(production, total_energy, 1);
linear_pred = polyval(linear_model, production);

quad_model = polyfit(production, total_energy, 2);
quad_pred = polyval(quad_model, production);

SS_res_linear = sum((total_energy - linear_pred).^2);
SS_res_quad = sum((total_energy - quad_pred).^2);
SS_tot = sum((total_energy - mean(total_energy)).^2);

R2_linear = 1 - (SS_res_linear / SS_tot);
R2_quad = 1 - (SS_res_quad / SS_tot);

RMSE_linear = sqrt(mean((total_energy - linear_pred).^2));
RMSE_quad = sqrt(mean((total_energy - quad_pred).^2));

fprintf('Linear R2 = %f\n', R2_linear);
fprintf('Quadratic R2 = %f\n', R2_quad);
fprintf('Linear RMSE = %f\n', RMSE_linear);
fprintf('Quadratic RMSE = %f\n', RMSE_quad);

figure;
scatter(production, total_energy, 'filled');
hold on;
plot(production, linear_pred, 'LineWidth', 2);
plot(production, quad_pred, 'LineWidth', 2);
xlabel('Production');
ylabel('Total Energy (kWh)');
title('Production vs Total Energy');
legend('Actual Data', 'Linear Fit', 'Quadratic Fit');
grid on;
\end{lstlisting}

\section{DG Reduction Scenario Code}
\begin{lstlisting}[language=Matlab]
dg_reduction = [0.10 0.30 0.50];

for i = 1:length(dg_reduction)
    reduced_dg = dg_energy .* (1 - dg_reduction(i));
    new_cost = grid_energy .* grid_tariff + reduced_dg .* dg_tariff;
    cost_saving = sum(cost) - sum(new_cost);
    co2_reduction = sum(dg_energy - reduced_dg) * emission_factor;
    
    fprintf('DG Reduction = %d%%\n', dg_reduction(i)*100);
    fprintf('Cost Saving = Rs. %.2f\n', cost_saving);
    fprintf('CO2 Reduction = %.2f kg\n', co2_reduction);
end
\end{lstlisting}

\section{Anomaly Detection Code}
\begin{lstlisting}[language=Matlab]
window = 7;
rolling_mean = movmean(total_energy, window);
rolling_std = movstd(total_energy, window);

anomaly = abs(total_energy - rolling_mean) > 2 * rolling_std;

figure;
plot(total_energy, 'LineWidth', 2);
hold on;
plot(find(anomaly), total_energy(anomaly), 'ro', 'MarkerSize', 8);
xlabel('Day');
ylabel('Total Energy (kWh)');
title('Energy Anomaly Detection');
legend('Energy Consumption', 'Anomaly');
grid on;
\end{lstlisting}

\section{Load Scheduling Logic}
\begin{lstlisting}[language=Matlab]
safe_grid_limit = 9740;
forecast_energy = 10255;

if forecast_energy > safe_grid_limit
    dg_load = forecast_energy - safe_grid_limit;
    
    unoptimized_cost = safe_grid_limit * grid_tariff + dg_load * dg_tariff;
    optimized_cost = forecast_energy * grid_tariff;
    
    savings = unoptimized_cost - optimized_cost;
    co2_avoided = dg_load * emission_factor;
    
    fprintf('DG Risk Detected\n');
    fprintf('Load to shift = %.2f kWh\n', dg_load);
    fprintf('Cost Saving = Rs. %.2f\n', savings);
    fprintf('CO2 Avoided = %.2f kg\n', co2_avoided);
else
    fprintf('Safe Load Condition. DG not required.\n');
end
\end{lstlisting}
